\documentclass[12pt]{article}%
\input setup

\begin{document}

\title{Mahler 2 Notes}
\author{David Hurwitz}
\date{}
\maketitle

\noindent
These notes are transcribed from the CD booklet that comes with the Micheal Gielen recording of this piece on Hanssler Classic.

\section{Mahler: Symphony No. 2}

At the time that it was completed in 1896, Mahler's »Resurrection« Symphony was among the largest compositions in the history of music, and so it remains to this very day. A revolutionary work in every respect, it's easy to imagine the shock that contemporary audiences must have felt at seeing so many singers and players participating in a work called a
»symphony.« After all, orchestrally speaking, the German symphonic tradition is nothing if not conservative. Both Bruckner and Brahms had remained faithful to the classical concept of the symphony orchestra as essentially a large string band opposed by a handful of woodwind, brass, and very limited percussion (usually timpani alone, maybe triangle and, more rarely, cymbals). Even Richard Strauss at his most extravagant never challenged this traditional instrumental hierarchy.

Mahler's symphonies changed all that. He was the first great symphonist to demonstrate that virtuoso orchestration could exist side by side with genuine, symphonic development. He was able to do this because his style was, from the first, essentially polyphonic: that is »horizon-tal, « concerned with the simultaneous presentation and development of themes and motives. Although his huge orchestra unleashes the most powerful sounds ever heard in a symphonic context, Mahler's fundamental goal is always clarity of texture first and foremost. To this end he has selected his instruments so that each section - winds, brass, strings and even percussion - possesses equal strength, as well as the ability to generate a completely independent mass of harmony and tone color within itself. Mahler's orchestra is really several orches-tras, each sometimes operating within the whole with an amazing degree of autonomy.
\newpage
The Second Symphony is scored for the following instrumental and vocal forces:

\begin{itemize}
\item 4 flutes (all four doubling piccolos)
\item 4 oboes (two doubling English horns)
\item 3 clarinets in B-flat, A and C (third doubling bass-clarinet)
\item 2 clarinets in E-flat (second doubling fourth clarinet in B-flat)
\item 4 bassoons (fourth doubling contrabassoon)
\item 10 horns (four offstage, later onstage)
\item 8 - 10 trumpets (six onstage, four offstage)
\item 4 trombones tuba
\item 7 percussionists in charge of:
\begin{itemize}
\item 2 sets of timpani (three drums each set; maximum of three players)
\item 1 kettledrum offstage
\item crash cymbals (1 pair offstage)
\item suspended cymbals several snare drums
\item 2 bass drums (1 offstage)
\item 2 triangles (1 offstage)
\item 2 tam-tams (high and low pitched)
\item glockenspiel
\item rute (bundle of sticks struck against the bass drum shell)
\item three low bells (steel bars of deep, differentiated but non-specific pitch)
\end{itemize}

\item 2 harps (several players per part)
\item Strings as numerous as possible: first and second violins violas cellos
\item basses (several with low C extension)
\item organ
\item soprano and alto soloists mixed chorus
\end{itemize}

\noindent
{\it Note}: In the following discussion, all timings listed are in minutes and seconds from the beginning of the movement in question.

The very opening of the symphony demonstrates Mahler's instrumental layering technique with the utmost clarity. Cellos and basses in unison state the long principal theme, and when the woodwind begin their independent melody, the lower strings continue with ther own music. Even when the clarinets, horns, and violins add weight to the upper melodic line, the cellos and basses keep right on, perfectly
audible beneath. Suddenly, the entire brass section enters as the whole orchestra comes together in unison for the first climax, then the players immediately divide again. In fact, the cellos and basses continue with their own mu-sic, virtually uninterrupted, for the first five minutes of the symphony.

Mahler's ideas thus may range from tiny frag-ments, to huge, seemingly endlessly unfolding melodies, while his constantly shifting orchestration presents the ear with a veritable kaleidoscope of sound. It's also one of Mahler's great achievements that many of his most fascinating and revolutionary innovations occur at quiet dynamic levels. Listen carefully, for example, at 4:45 into the first movement, to the slow trudge of basses backed by quiet stokes on the large tam-tam (gong). This is something new, and it's one of Mahler's most characteristic »funeral« effects - a cold, dark sound that will make you shiver. The basses finally fade away, leaving only the lowest notes of the two harps. Such a simple effect, but where have we ever heard anything like it before?

The next passage, which begins at 5:50, is an orchestral miracle, and a classic instance of Mahler's use of a huge orchestra at the lowest dynamic level to create the most delicate threads of sound. Violins softly ascend, as the horns begin an independent duet, only to be interrupted after a few bars by a chord on the harps. A soft call on the lower horns introduces a descending line first on oboes, then clarinets.
The trumpets repeat the horn call, then a new tune on solo English horn yields to the oboes and violins, with a countermelody in the bass clarinet. Next, two clarinets have yet another tune in duet, and the harp has the bass-cla-rinet's countermelody, followed once again by the horns, violins, and finally the cellos (taking up the English horn tune). And all of this takes only about a minute! The colors are constantly changing, always in perfect instrumental bal-ance, without ever breaking the long, seemingly endless lines of melody. This is what Mahler's scoring is all about.

The first movement is, basically, a huge and stormy funeral chant in which softer, lyrical elements constantly yield to the stern march rhythms of the cellos and basses. The conflict between these two kinds of music creates the movement's dramatic tension, and the march elements are the clear winners, though the ghostly and dark coda (beginning at 18:48: more trudging basses and funeral tam-tam strokes) proves the victory to be hollow indeed.
The second movement can be seen as the opposite of the first. Again, two types of music are contrasted: the lyrical opening theme, and a more agitated section in triplet rhythm (1:15).
Their charm and warmth are deceptive, as both themes are polyphonic: the opening theme returns accompanied by a gorgeous new melody in the cellos, while the more agitated music sounds at times like the famous fugal opening of the Scherzo of Beethoven's Ninth Symphony. Unlike the first movement, however, it's the lyrical element that dominates, nowhere more so than when it comes back on plucked strings, accented by harps and piccolo.
The third movement is an expanded version of one of Mahler's most humorous and ironic songs, »St. Anthony of Padua's Sermon to the Fish. « It's not necessary to know what the song is about to follow the music in its watery, tipsy, slightly mocking course. A sharp slam on the timpani sets the music going, propelled by the
»boom click click« rhythm on the bass drum and rute. At 4:33, Mahler unveils another one of his orchestral masterstokes - a sentimental, Romantic tune played by a quartet of trumpets.
It's moments like this that justify the presence in Mahler's symphonies of so many of the same type of instrument. The result is simply magical.
Even a musical conservative such as Brahms, no great fan of Mahler the composer, loved this music.

The movement climaxes with what Mahler calls
»a cry of despair,« and it certainly sounds like one. But what's even more important is that this huge eruption embodies one of Mahler's favorite dramatic techniques: the triumphant climax that somehow »goes wrong, « becoming terrifying or tragic in the process. We will encounter something similar more than once in the finale.

The last three movements of the symphony are all played without a pause. A quiet stroke on the tam-tam ends the third movement, out of which the contralto soloist begins the brief song, »Urlicht« (»Primal Light«), which constitutes the fourth. While not the first time that voices had been used in a symphony, Mahler was the very first composer to incorporate whole songs into a symphonic scheme. As with all texted music, the best way to understand it is simply to follow the words as they are sung.
They express a simple message: the desire of the soul to be accepted into heaven by a loving God. Pay special attention to the song's middle section which begins at 2:45 (»I came upon a broad path...«), and to the lovely writing for glockenspiel (small bells), clarinet, harp, solo violins, and two piccolos.

Mahler's answer to the »questions« posed by the widely differing emotional states of the first three movements is much more complicated than the simple message of hope offered by a mere song, however solemn and lovely. The orchestra rejects the fourth movement, and the finale begins with a crash, and the »cry of despair« which we have already heard just a few minutes previously. You will notice, though, that the explosion of sound isn't exactly the same as before. The trumpets have an actual tune on top of the »cry of despair,« and this tune is the key to most of the music that you will hear in the half hour to come. The trumpet theme has two parts: a musical proclamation of two three-note phrases (:05), followed by a simple fanfare (:10). As the eruption calms down, a third, ascending theme appears quietly in the horns like a vision of paradise seen from afar (:35): this is the theme of »resurrec-tion,« the movement's eventual goal.

Now comes a series of ten distinct episodes, a giant set of variations on this titanic opening, as well as themes from earlier in the symphony, each variation revealing a different stage on the journey to the symphony's transfigured ending.
These episodes are very clearly demarcated, and so magnificently planned and characterized that they deserve consideration in some detail.

{\bf First Episode} (1:50): Offstage horns call from the distance. The orchestral responds with a brief passage featuring oboe, trumpets and onstage horns that descends via mysterious trilling chords on the strings and harps to a quiet funeral march, complete with those (by now famous) tam-tam strokes.

{\bf Second Episode} (3:19): Woodwinds play a tune that may sound familiar: it actually comes from the climax of the first movement (if you want to hear for yourself, go back to 12:15 in that movement and listen to the horns). A solo trombone, then solo trumpet take over the tune, and suddenly the opening fanfare returns on the horns, only purged of all urgency and suffused with quiet light. The offstage horns then have a more melodic version of their initial horn call leading back, once again, to the funeral march which ended the First Episode.

{\bf Third Episode} (5:45): An exercise in musical anxiety, lead off by a pleading English horn, shrieking clarinets, and jagged string tremolos.
The musical material is all new.

{\bf Fourth Episode} (6:58): This is an almost exact repetition of the Second Episode, except that it's entirely rescored. Where it once sounded timid, now it appears solemn and confident.
The full brass section plays a richly harmonized chorale version of the tune from Episode Two (which in turn came from the first movement).
This builds to a huge climax, capped by a varied return of the theme on the trumpets that opened the movement. Only now it's played by horns and trombones and in reverse phrase or-der: the fanfare comes first, alternating with statements of the three note phrases. The volume is huge, but so carefully balanced is Mahler's scoring that if you listen, you can actually hear the harps thrumming underneath it all. This leads, via the same music as in the Second Episode, back once again to the funeral march, now at its most threatening.

{\bf Fifth Episode} (10:15): Two huge solo percussion crescendos (snare drums, tam-tams, bass drum, and timpani) separate full brass section statements of those two, three-note phrases from the movement's opening. Mahler then repeats these in different positions over music of real violence that eventually becomes a full-blown, parade ground march. Strings and timpani begin, and then a fully developed marching version of this opening tune comes back on the trumpets while the strings continue beneath (a technique we should now recognize as typically Mahlerian). This march goes through several episodes, leading to a dramatic statement of the first movement tune from Episodes Two and Four. In fact, for a moment we seem to be back in that movement - the approach to the climax is exactly the same. A pause, and with a huge crash (tam-tam, cymbals, bass drum and three timpanists) the entire episode simply self-destructs.

{\bf Sixth Episode} (14:43): The music is that of the Third Episode, begun on trombone, and becoming a lyrical cello theme that plays against offstage percussion with added trumpet fanfares.
This gathers itself up into a huge climax that culminates in a return to the »cry of despair« and the very opening of the movement. The noise stops, and we find ourselves essentially back where we started, after all of this effort. What happens next, the music seems to ask?

{\bf Seventh Episode} (18:30): Offstage horns, trumpet and timpani over a quiet bass drum roll engage in a dialog with flute and piccolo in the orchestra. Mysterious fanfares, drum rolls and birdsong fill the air. Something special is about to happen. Something wonderful.

{\bf Eighth Episode} (21:07): The chorus, unaccom-panied, sings softly of the promise of resurrec-tion. »Arise, yes, you will arise.« The words are from Klopstock's »Resurrection Ode, « as adapted and added to by Mahler himself. The music is, of course, a very slow, solemn version of the opening trumpet tune in its Fifth Episode, march variation. As the chorus sings, a solo soprano voice floats free of the vocal mass, and the orchestra answers with the »heavenly« version of the opening fanfare, as heard way back in the Second Episode, combined with the »resurrection theme«. Mahler's scoring has never been more amazing than here, for although the music seems weightless and lit with an inner radi-ance, he's actually using the whole orchestra, including a full on-stage brass section of ten horns, six trumpets, four trombones and tuba!
The next verse of the Ode is set to the same music, and leads once again to the orchestral interlude which this time achieves a full close.

{\bf Ninth Episode} (27:50): The music of the Third and Sixth Episodes becomes a contralto solo, and now we understand the reason for the music's anxiety. »Oh Believe!« she sings, »You have not lived in vain. What you have striven for will be yours.« She's answered in turn by the solo soprano, and then the male chorus, which sings a variation of the music of the first choral entrance. »Prepare yourself to live!« the singers shout, answered by the contralto, rising out of the mass of voices as the soprano did twice previously.

{\bf Tenth Episode} (31:01): The two soloists sing ecstatically of the promise of eternal life. Their song is taken up by the chorus, combining the
»resurrection« theme with the music of their initial entrance, triumphantly accompanied by the full orchestra and organ. At their final cadence (»What you have struggled will bear you to God!«), the »resurrection« theme ascends into the heavens, accompanied by the pealing of bells. A final moment of quiet, like a deep sigh of contentment, leads to the final bars in which bells, tam-tams, organ and the full orchestra proclaim the joy of the Resurrection of mankind, and of the Mahler's conception of a merciful, loving, and above all non-sectarian Creator.

\end{document}
